%======================================================================
\section{Conclusions}
\label{sec:conclusions}
%======================================================================

A latent field, built only from operators that are shared by all sites
and read a bounded neighbourhood, gives a learned surrogate of extensive
chaos that can be run on domains it was never trained on, and the
models' own Lyapunov spectra show that it reproduces the density of
chaos there. In one dimension the transferred Kuramoto--Sivashinsky model
stays within $0.2$--$1.2\%$ of the true Kaplan--Yorke dimension up to four
times the training length, and its spectrum, not only its dimension,
follows the truth's. In two dimensions the architecture alone is not
enough: what makes the learned rule transfer is the data it is trained
on. Trained on a small periodic box it misses the dimension of larger
domains by up to $38\%$; trained on windows of one large simulation, with
the boundary of each window driven by the simulation, it is within
$0.6$--$2\%$ in 2D Kuramoto--Sivashinsky and $1$--$9\%$ in Kolmogorov flow at
the two stronger drags; of six training variants at the strongest drag,
the one chosen without any reference to the truth has the lowest mean
error. The weakest drag remains harder.

The failures are as instructive as the successes. Rollouts that are too
short leave the free-running model room to settle into spurious regimes,
and squared-error rollouts that are too long damp its growth rates, a
bias that follows from the structure of the loss; free-rollout
statistics, compared with the training data, catch the first problem
without any truth, and the dimension comparison exposes the second.

The same locality carries over to data assimilation. With a per-patch
decoder the latent observation operator is exactly local, so a per-site
ensemble filter in latent space is exact, needs only a fixed number of
forward passes for its Jacobian, and uses memory that grows no faster
than the area. With 64 members it keeps $90\%$ of its skill from $44^2$ to
$440^2$, while an unlocalized filter fails whenever the attractor has more
unstable directions than the ensemble has members; in one dimension a
localized filter needs 16 members at every size, and an unlocalized one
needs twice as many with every doubling of the domain.

The rule we take from this work is short: measure the interaction range
of the dynamics, give the model a reach that covers it, train on windows
of a domain large enough to contain the local statistics of the target,
and check the result against the Lyapunov spectrum of the truth at more
than one size. Where those conditions hold, a model trained once is a
model for every domain size, and a localized filter built on it does not
need a larger ensemble as the domain grows.
